% !TeX root = ../papers/nc-p3.tex
\section{平方零の三重点と長さ六の平面切断}
\label{sec:finite-marked-planar-triple}
\begin{theorem}\label{thm:finite-marked-planar-triple}
共通の印付き切断で階数1軌跡が非Gorensteinの長さ3なら、その四平面は
\[
 E_D=\left\{\begin{pmatrix}a&b^t\\ b&tD\end{pmatrix}:
       a,t\in k,\ b\in k^2\right\},\qquad D\in\mathrm{GL}_2,\quad D\ne D^t
\]
に左右基底変更で一致する。$D$ の同一視はスカラー倍と合同で尽きる。
印が曲線状または被約なら、この模型でも $T$ が反対位置の対称性を与える。
従って接次元1以下の印では長さ3の全ての階数1切断がこの対称性判定に含まれる。
\end{theorem}
\begin{proof}[模型と二つの平面]
非Gorensteinの長さ3の環は一つの点に支持され、その極大イデアルは二次元で平方零である。
支持行列を $M=E_{11}$ とする。Segreの接行列は右下二行二列が零である。
二つの接ベクトルの左下二座標も右上二座標も各々独立である。
一方で従属なら線型結合で $M$ と共通像または共通余像を持つ接方向があり、
その直線全体が $E$ 中の階数1行列になる。従って基底変更で双方を同じ $b$ とできる。
第四の行列の右下 $D$ は可逆である。さもなければその左核の元は $E$ 全体を階数1以下に縮約する。
この可逆性による行・列加算で第四行列の上段と左列を消せる。接二行列には $M$ の倍を
加えるだけである。$a=1$ で階数1の式は $tD=bb^t$。
$D$ が対称なら、$t$ を消去した二本の二次式は二変数平方零の三独立二次式を張れず矛盾する。
非対称なら $t=0$, $bb^t=0$ を得て、実際に平方零環になる。$a=0$ に他の階数一点はない。
この支持を保つ左右変換はそれぞれ $ke_1$ を保つ。二つの下側二行二列の変換は
共通の接変数 $b$ を対応させるので互いに比例する。ゆえに $D$ の変換はスカラー倍と合同であり逆も明らかである。

中央文字は $E_D$ の座標 $(a,b,t)$、位置1の座標は印の補座標と射影左方向 $l=(l_0,l')$ とする。
$T$ で同定した位置0,2の三平面内の射影平面を $P=\{t=0\}$、位置1の平面を $L=\{l_0=0\}$ とする。
$R$ を辺 $01$ の六関係、$R'$ を辺 $12$ を反転して $T$ で移した六関係とする。
両空間は $l^t E_D$ の三成分を共有する。実際これらは $p_2$ と $p_0$ で縮約した後に位置3で縮約した
もので、共通切断の同定そのものである。$P\times L$ では三式は $(l'^t b,0,0)$ に制限される。
$R$ の全制限の独立数は4以下で、その基点が有限なのでちょうど4である。
そこで $Y=X_{01}\cap(P\times L)$ は四双一次式の完全交叉で長さ6である。
その $(1,1)$ の制限像は次元5、位置1だけの二次式の制限は $H^0(\mathcal O_Y)$ への
自明化後の同型である。前者は消える四切断が定義式で尽きること、後者は四式のKoszul分解を
$(0,2)$ でひねった負の全項のコホモロジーが消えることによる。

\end{proof}
\begin{proof}[誤差と三行三列の余因子]
平方零環の相対割線は全て $P$ にあり、左核は $L$ に対応する方向にある。
実際、階数1の切断に含まれる相対次数2の部分の線型包は $P$ に入る。
その鉛筆は剰余で $M$ と一つの非零接行列に張られる。相対的にも $M$ を含む：
平方零環の基底を $1,x,y$ とすると、相対直線イデアルの生成元
$u+vx+wy$ は $v,w$ の一つが単元であり、$x,y$ の乗法でその直線を保つため $u=0$ である
（例えば $v$ が単元なら $y$ による積 $uy$ が生成元の倍であることからその倍数が零になる）。
従って長さ2の商は定数点を含む。$M$ を消す左核は $l_0=0$ を満たす。

$X=X_{01}$ 上で $R'$ を評価する。印以外の局所因子で差 $Z=f_0 V$ があれば、
$\operatorname{ann}(f_0)$ を法として、前節までの消去で独立な二元 $A,V$ が相対割線を張る。
したがって任意の評価 $f$ について $t(x_0) f=l_0(x_1) f=0$ である。
ここで位置0の $t$ は $T$ を通した座標、$x_0$ は適宜 $T x_0$ と読む。
印の環を $k[z]/(z^m)$ とし、$Z=z^s V$ とする。$z^{m-s-1}$ を法として同じ結論を
$l/z,r/z$ に適用すると $l_0(x_1) f=0$ で、$t(x_0) f$ は少なくともsocleに入る。
$m-s=1$ の場合も $l_0$ が印で零なので同様、被約の印では差は零である。
実は $t(x_0) f=0$ である。これは印以外で零、印でsocleに入り、
$\tau_X(t(x_0) f U_1)=0$ は隣接座標積の自己直交性である。$U_1$ は印で単元を含むので
完全跡により零となる。ゆえに全評価は $C\cap\operatorname{ann}(I_Y)$ に入る。
この空間は次元1である。実際 $C=C^\perp$ よりその次元は $C+I_Y$ の余次元、
すなわち $6-5$ に等しい。

非零評価 $f$ があると仮定する。その支持は補題\ref{lem:finite-adjacent-word-support}より
長さ6以上で、$Y$ に含まれるから $Y$ に一致する。$H=R\cap R'$ は五次元である。
各局所環で二つの偶座標の差の成分が張るイデアルは $(f)$ と等しい。
実際、第二辺の位置1での評価行列は $X$ 上で階数3で、核直線は $x_2$ である。
単元な三次小行列で消去すれば、その六式の $T x_0$ での値が張るイデアルは
規格化した $T x_0,x_2$ の座標差が張るものに等しい。これが共通四平面の $Z$ の成分のイデアルである。
一方で全評価は $f$ の体スカラー倍である。

共有三式の制限が一本なので、$H$ は $P\times L$ 上の三本以下の双一次式を与える。
さらに一本で長さ6の有限完全交叉を切るのでちょうど三本で、その三行三列の係数行列を
$J(x_1)$ とする。これは $L$ の一般点で可逆である。そうでなければ三式の零点が平面 $L$
を覆って次元2以上となり、さらに一本のSegre超平面では有限にならない。
印以外の $Y$ 上では $A,V$ が $P$ に入り $H(x_1)$ の独立な二核元となる。
実際 $H$ は両方のグラフを消し、差の共通生成元で割れば $V$ を $\operatorname{ann} f$ を法として消す。
印では同じことが局所長さを1だけ減らせば成り立つ。上で $z^{m-s-1}$ を法として得た
相対割線全体が $P$ に入るからである。従って $J$ の全二次小行列式は、$Y$ の長さ5以上の部分で消える。
位置1の二次式が $Y$ の長さ6を分離するので、全余因子は同一の二次式の定数倍である。
一般可逆性から $\operatorname{adj}J=qK$ の定数行列 $K$ は可逆である。
随伴行列の積の式で $J$ 自身が一本の一次式の $K^{-1}$ 倍となり、その直線上で階数0に落ちる。
第四の一式を加えてもその上で有限交叉とならず矛盾する。従って全評価は零である。
\end{proof}
